\subsection{Technical Foundations and System Design Rationale}
\subsubsection{Component-Based Interface Design and Cognitive Load}
The practical software framework for a learning system must be architected to respect the human cognitive architecture, specifically the capacity and duration limits of working memory. To support low-latency and highly interactive applications, the frontend architecture utilizes component-based UI rendering, which facilitates the "segmenting effect" by breaking complex, high-element interactivity lessons into manageable, learner-controlled segments. This approach ensures that learners are not overwhelmed by a single continuous unit of information, allowing working memory resources to recover between processing episodes. Furthermore, a strongly-typed backend architecture is necessary to manage the "state-heavy" nature of individualized learner modeling, where the system must track millions of unique parameters—such as item-specific Easiness Factors (EF), Memory Stability (S), and Memory Retrievability (R)—across high-dimensional matrices. By minimizing extraneous load through a simplified, responsive user interface, the framework ensures that the learner’s limited cognitive resources are focused entirely on the germane task of schema construction\cite{Cog2019Sweller, mayer2003, studiomg_cognitive_ai}.

\subsubsection{Integration of the Adaptive Loop}
The system operates through a continuous "use, verify, and correct" feedback loop that synchronizes the web client, vector database, LLM, and SM-2 based algorithms. The loop begins when the web client presents a learning item; once the student provide a self-assessed quality grade (0–5), this "retrievability" data is sent back to the algorithm to verify the accuracy of the predicted interval. If the grade is low, indicating forgetting, the system immediately "corrects" the multidimensional matrices by reducing the stability increase factors and resetting the repetition schedule. Simultaneously, the LLM utilizes RAG to pull "supportive information" from the vector database, identifying semantically related knowledge nodes through high-dimensional search to provide "just-in-time" procedural hints. This integration ensures that the algorithm "learns the learner" in real-time, shifting from a fixed, heuristic model to a personalized forgetting curve that adapts to individual memory decay rates\cite{wozniak1998sm2, pokrywka2023lstm_srs}.